\section{What is built once, and what every query pays}
\label{sec:shared}
\Needspace{12\baselineskip}
\subsection{The public Exa-style starting point}
Exa's public description combines Matryoshka prefixes, binary document codes, float-query
scoring through small lookup tables, similarity-based clustering, filtering, and final
reranking \cite{exa}. Its published example uses particular dimensions and cluster counts;
our formulas replace those choices with parameters. The pseudocode here is an explicit
reference design, not a claim to reproduce Exa's private implementation or production latency.

\Needspace{12\baselineskip}
\subsection{Preprocessing and the common index}
A document is encoded before queries arrive. Its code and document ID are placed in its
cluster's contiguous row range. The baseline ownership convention is simple: each document
has one packed code and one ID. Cluster offsets identify the ranges. Bitplanes or hash-key
postings are additional structures described later.

\par\noindent\begin{minipage}{\linewidth}
\begin{lstlisting}[caption={Common preprocessing; executed before serving queries.}]
BuildCommonIndex(documents, model, d, routing_configuration):
    vectors = encode documents with model
    routing = train or choose the cluster assignment rule
    for each document x:
        b = signs of the first d coordinates of vectors[x]
        j = assign x to one cluster
        append (document_id[x], packed(b)) to cluster j
    concatenate cluster rows and record C + 1 offsets
    save centroids if centroid routing is used
    save filter metadata and optional final vectors
    return common index
\end{lstlisting}
\end{minipage}\par

Offline encoding, clustering and sorting are build costs. They are not charged as query
latency. Their outputs \emph{are} charged as index storage. Build-time temporary arrays can be
large; the serving-RAM formulas below deliberately do not describe the construction peak.

For centroid routing, define
\begin{equation}
\Mbase=N s_{\rm code}+N b_{\rm id}+(C+1)b_o+C d_r b_c+M_{\rm filter}+M_{\rm extra}.
\label{eq:base}
\end{equation}
$M_{\rm extra}$ includes explicitly chosen common metadata and optional stored final vectors.
For example, fully resident reranking vectors contribute $N d_f b_f$. If they are on disk,
they still consume storage, but only the resident portion belongs in serving RAM. For
sign-prefix routing, omit the centroid term. Do not silently count a second ID array in one
method and omit the only ID array in another.

\begin{examplebox}{Stored bytes versus touched bytes}
Suppose an illustrative database has $N=10^9$ documents and $d=256$ bits. Packed code payload
without extra padding is $10^9\cdot256/8=32$ GB. With eight-byte IDs, the codes and IDs together
need 40 GB before centroids or metadata. Here GB means $10^9$ bytes. These are example choices
for $N,d,b_{\rm id}$, not constants in Equation~\eqref{eq:base}.
\end{examplebox}

\Needspace{12\baselineskip}
\subsection{Step 1: encode and prepare the query}
\par\noindent\begin{minipage}{\linewidth}
\begin{lstlisting}[caption={Query preparation shared by all methods.}]
PrepareQuery(text, model, D, d, d_r):
    full_query = model.encode(text, D)
    q = prepare prefix full_query[0:d] for the chosen score
    routing_query = prepare prefix full_query[0:d_r]
    t[i] = +1 if q[i] >= 0 else -1
    a[i] = abs(q[i])
    return full_query, q, routing_query, t, a
\end{lstlisting}
\end{minipage}\par
The model's tokenization, neural-network computation and output processing cost
$T_{\rm emb}(\ell_q,D)$ for query length $\ell_q$. This is a measured function, not a universal
linear formula in $D$. Truncation, optional normalization, signs and magnitudes add
$T_{\rm prep}=O(D+d+d_r)$ elementary work. Their coefficient can be measured separately.
A norm or layer-normalization rule must be the same rule used when preparing the documents.

Query-vector storage is $O(D+d+d_r)$ floats, with possible views instead of copies. Record the
chosen ownership. We denote the actual retained query buffers by $Q_{\rm vec}$ and the
embedding model's own scratch peak by $Q_{\rm emb}$.

\Needspace{12\baselineskip}
\subsection{Step 2: select clusters}
With a flat centroid router, score all $C$ centroids, then retain the best $P$. This is distinct
from scoring binary documents: centroids are normally floats.
\par\noindent\begin{minipage}{\linewidth}
\begin{lstlisting}[caption={A simple flat centroid router with bounded selection memory.}]
Route(q_route, centroids, P):
    nearest = heap with capacity P
    for j in all C clusters:
        route_score = dot(q_route, centroid[j])
        offer (route_score, j) to nearest
    return the P selected cluster IDs
\end{lstlisting}
\end{minipage}\par
Let $c_{\rm fma}$ price one centroid coordinate's multiply-add contribution, and let
$H_r$ be the number of heap comparisons or equivalent selection operations. Then
\begin{equation}
T_{\rm route}=C d_r c_{\rm fma}+H_r c_{\rm sel},\qquad
H_r=O(C\log(P+1)).
\label{eq:route}
\end{equation}
A streaming heap needs $O(P)$ records; materializing all centroid scores needs $O(C)$ scores
instead. Either is valid, but their RAM formulas differ. Hierarchical or approximate routing
has its own measured visited-centroid count; replace $C$ by that count only when the algorithm
actually avoids the other centroids. Its routing recall also needs evaluation.

A sign-prefix router uses $r$ fixed sign coordinates and at most $2^r$ addresses. It is not the
same as similarity-based centroid clustering. Assigning a document is cheap, but ranking many
possible query buckets is not free. One reference router scores all $C$ bucket keys using
$\ceil{r/g}$ lookup terms each and selects $P$, giving
\[
T_{\rm route,sign}=T_{\rm table,r}+C\ceil{r/g}c_{\rm lut}+H_r c_{\rm sel}.
\]
A direct or multi-probe address generator can replace that router, but must count its own
key-generation and queue work. When every cluster is probed, routing can simply return all
cluster IDs and skip distance-based selection.

\Needspace{12\baselineskip}
\subsection{Cluster count does not determine the selected document count}
Exactly,
\begin{equation}
\Nsel(q)=\sum_{j\in\mathcal P(q)}n_j.
\label{eq:selected}
\end{equation}
Only under balanced clusters is this approximately $NP/C$. An optional load factor
$\lambda(q)=\Nsel/(NP/C)$ summarizes deviation for a \emph{given query}; it is not automatically
constant as $P$ changes. For bitplanes, the individual $n_j$ matter because wider clusters
make each node more expensive. A mean load factor does not identify the largest cluster or
peak frontier memory.

\Needspace{12\baselineskip}
\subsection{Step 3: apply optional filters}
A filter can return eligible row IDs or a bitmap. Denote its time by $T_{\rm filter}$ and
its scratch by $Q_{\rm filter}$. For a bitmap intersection over the probed clusters, a reference
work count is $\sum_{j\in\mathcal P}W_j$ word visits per input mask. For sorted ID lists, cost
depends on list lengths and the intersection method instead.

After filtering, $e_j$ rows are eligible in cluster $j$, so a scan scores
$\Nelig=\sum_j e_j$ documents. A dense bitplane still has physical width $W_j$. A filter that
removes most rows therefore helps scoring immediately, but does not automatically make each
bitplane intersection short.

The reference set for recall must obey the \emph{same filter}. Otherwise a correct filtered
search is unfairly penalized for not returning ineligible documents.

\Needspace{12\baselineskip}
\subsection{Step 4: build the float-query score tables}
Split the $d$ document signs into groups of $g$. A group has at most $2^g$ sign patterns.
For each query group, precompute its dot product with every possible pattern.
\par\noindent\begin{minipage}{\linewidth}
\begin{lstlisting}[caption={The score table. Padding contributes zero in the last group.}]
BuildScoreTable(q, g):
    for group in ceil(len(q) / g) groups:
        for pattern in all g-bit patterns:
            table[group, pattern] = 0
            for coordinate in this group:
                add q[coordinate] * sign(pattern, coordinate)
    return table

ScoreDocument(packed_code, table):
    score = 0
    for group in all groups:
        pattern = extract this group's bits from packed_code
        score += table[group, pattern]
    return score
\end{lstlisting}
\end{minipage}\par
For this straightforward build,
\begin{align}
G&=\ceil{d/g}, & U_{\rm build}&\leq G2^g g,\\
T_{\rm table}&=U_{\rm build}c_{\rm build}, & Q_{\rm table}&=G2^g b_s,\\
T_{\rm score}(F)&=F G c_{\rm lut}.&&
\label{eq:lut}
\end{align}
$c_{\rm lut}$ includes extracting the pattern, the lookup and accumulation in the chosen
kernel. Faster table construction is possible, but it must replace $U_{\rm build}$ explicitly.
Larger $g$ reduces the number of lookup terms and increases the table exponentially. Register
capacity and cache behavior therefore matter.

\begin{examplebox}{Two query values and a four-entry table}
For the first two values $(0.7,-0.4)$, the contributions for patterns
\texttt{00}, \texttt{01}, \texttt{10}, \texttt{11} are respectively
$-0.3,-1.1,1.1,0.3$. With $g=2$, document \texttt{1011} uses the first-group entry
$1.1$ and the second-group entry $0.1$, giving $1.2$.
\end{examplebox}

Plain XOR followed by a population count computes \emph{unweighted} Hamming distance. It
matches Equation~\eqref{eq:score} only when all query magnitudes are equal up to a common
scale. Replacing weighted scoring by that operation changes the task.

\Needspace{12\baselineskip}
\subsection{Step 5: retain candidates and optionally rerank}
A top-$R$ heap stores the worst retained candidate at its root. Every scored document needs
at least a comparison with that root once the heap is full. Let $H_s(F,R)$ be the actual
selection-operation count. A simple worst-case bound is $O(F\log(R+1))$, not zero.
\begin{equation}
T_{\rm select}=H_s(F,R)c_{\rm heap},\qquad
Q_{\rm results}\approx R(b_{\rm id}+b_s)+Q_{\rm heap,overhead}.
\label{eq:selection}
\end{equation}
For an optional rerank of $R'$ returned candidates, $K\leq R'\leq R$, let
\begin{equation}
\Tfinish(R')=T_{\rm fetch}(R',d_f)+R'd_f c_{\rm fine}
              +H_f(R',K)c_{\rm fine,sel}+T_{\rm output}(K).
\label{eq:finish}
\end{equation}
$T_{\rm fetch}$ depends on whether vectors are already resident, on SSD, or remote. If $u_f$
vectors are fetched in a batch, the temporary float-vector buffer is $u_f d_f b_f$, not
necessarily $R'd_f b_f$. The candidate IDs must remain live until fetching and scoring finish.

\Needspace{12\baselineskip}
\subsection{Shared latency and peak RAM}
For sequential stages, define
\begin{equation}
\Tcommon=T_{\rm emb}+T_{\rm prep}+T_{\rm route}+T_{\rm filter}+T_{\rm table}.
\label{eq:common-time}
\end{equation}
The method-specific candidate work and Equation~\eqref{eq:finish} are added later. For an
index-only measurement, explicitly set the embedding and final-stage terms to zero.

Peak scratch RAM is a maximum of simultaneously live buffers, not a sum of every temporary
array ever allocated. If embedding, search and reranking do not overlap,
\begin{equation}
Q_{\rm peak}=\max\{Q_{\rm emb,stage},Q_{\rm search,stage},Q_{\rm finish,stage}\}.
\label{eq:peak}
\end{equation}
Each stage term includes buffers retained from previous stages. We use $\Qbase$ for common
buffers live during candidate search, including query vectors, score tables, filter masks,
routing output, and the candidate heap.

For $J$ concurrent queries sharing one resident index,
\begin{equation}
\begin{aligned}
M_{\rm serving}&=M_{\rm resident}+M_{\rm model}+M_{\rm runtime}
                 +\max_t\sum_{u\in\mathcal A(t)}Q_u(t)\\
&\leq M_{\rm resident}+M_{\rm model}+M_{\rm runtime}+JQ_{\rm peak,max}.
\end{aligned}
\label{eq:serving}
\end{equation}
$M_{\rm model}$ counts shared resident embedding-model weights when inference is local. An external embedding service contributes latency but no local model-weight allocation. $M_{\rm runtime}$ is a chosen reserve or measured runtime overhead. It is not another stored
copy of the index. Separate process replicas or device copies must be counted if the actual
system has them.

For separate CPU and GPU memories, apply the residency and peak formulas to each memory pool. Count additional device copies if they exist. A throughput speedup does not by itself establish a feasible memory layout.
